\documentclass{article}
\usepackage{amsmath,amssymb,amsthm}
\usepackage{algorithm}
\usepackage{algpseudocode}
\usepackage{enumitem}
\newtheorem{theorem}{Theorem}
\newtheorem{lemma}{Lemma}
\newtheorem{assumption}{Assumption}
\newtheorem{definition}{Definition}
\newcommand{\norm}[1]{\left\lVert#1\right\rVert}
\begin{document}

\section*{Algorithm Name}
\textbf{Trust‑Region Gradient Method (TR‑GM)} for convex and smooth objectives.

\section*{Mathematical Formulation}
Given a differentiable convex function $f:\mathbb{R}^n\to\mathbb{R}$ with $L$‑Lipschitz gradient, the $k$‑th iteration solves the subproblem
\[
\min_{p\in\mathbb{R}^n}\; m_k(p)\;\;:=\;\;f(x_k)+\nabla f(x_k)^{\top}p+\frac{1}{2}p^{\top}B_k p
\quad\text{s.t.}\quad \norm{p}\le \Delta_k,
\tag{1}
\]
where $B_k$ is a symmetric positive semidefinite matrix.  In the convex smooth setting we take the simple choice
\[
B_k = L I_n,
\tag{2}
\]
so that $m_k(p)$ is a quadratic upper model of $f$.

Let $p_k$ be an (possibly inexact) solution of (1).  Define the actual reduction
\[
\text{ared}_k \;:=\; f(x_k)-f(x_k+p_k),
\]
and the predicted reduction
\[
\text{pred}_k \;:=\; m_k(0)-m_k(p_k)
            \;=\; -\nabla f(x_k)^{\top}p_k-\frac12 p_k^{\top}B_k p_k .
\]
The ratio
\[
\rho_k \;:=\; \frac{\text{ared}_k}{\text{pred}_k}
\tag{3}
\]
governs acceptance of the trial step and the update of the trust‑region radius $\Delta_k$.

Parameters:
\begin{itemize}[nosep]
\item $\eta_1\in(0,1)$ – acceptance threshold (commonly $0.1$).
\item $0<\gamma_{\text{inc}}<\gamma_{\text{dec}}$ – radius expansion/shrink factors (e.g. $\gamma_{\text{inc}}=2$, $\gamma_{\text{dec}}=0.5$).
\item $\Delta_{\max}>0$ – maximal radius.
\item $\Delta_0>0$ – initial radius.
\end{itemize}

Update rules:
\[
x_{k+1}= 
\begin{cases}
x_k+p_k, & \text{if }\rho_k\ge \eta_1,\\[4pt]
x_k,    & \text{otherwise},
\end{cases}
\tag{4}
\]
\[
\Delta_{k+1}= 
\begin{cases}
\min\{\gamma_{\text{inc}}\Delta_k,\;\Delta_{\max}\}, & \text{if }\rho_k\ge \eta_1,\\[4pt]
\gamma_{\text{dec}}\Delta_k, & \text{if }\rho_k<\eta_1 .
\end{cases}
\tag{5}
\]

A convenient inexpensive approximate solution of (1) is the \emph{truncated gradient step}
\[
p_k^{\text{tg}} \;:=\; 
\begin{cases}
-\dfrac{1}{L}\nabla f(x_k), & \text{if }\norm{\nabla f(x_k)}\le L\Delta_k,\\[6pt]
-\Delta_k\,\dfrac{\nabla f(x_k)}{\norm{\nabla f(x_k)}}, & \text{otherwise}.
\end{cases}
\tag{6}
\]
Throughout the analysis we set $p_k:=p_k^{\text{tg}}$.

\section*{Pseudocode}
\begin{algorithm}[H]
\caption{Trust‑Region Gradient Method (TR‑GM)}\label{alg:trgm}
\begin{algorithmic}[1]
\Require Convex $L$‑smooth $f$, initial point $x_0$, parameters 
$\Delta_0,\Delta_{\max},\eta_1,\gamma_{\text{inc}},\gamma_{\text{dec}}$.
\State $k\gets0$, $\Delta_k\gets\Delta_0$.
\While{stopping criterion not satisfied}
    \State Compute gradient $g_k\gets\nabla f(x_k)$.
    \State \textbf{(Trial step)} Compute $p_k$ via (6).
    \State Compute actual reduction $\text{ared}_k = f(x_k)-f(x_k+p_k)$.
    \State Compute predicted reduction $\text{pred}_k = -g_k^{\top}p_k-\frac12 p_k^{\top}(L I)p_k$.
    \State $\rho_k\gets \dfrac{\text{ared}_k}{\text{pred}_k}$.
    \If{$\rho_k\ge \eta_1$}
        \State $x_{k+1}\gets x_k+p_k$ \Comment{Accept step}
        \State $\Delta_{k+1}\gets \min\{\gamma_{\text{inc}}\Delta_k,\;\Delta_{\max}\}$.
    \Else
        \State $x_{k+1}\gets x_k$ \Comment{Reject step}
        \State $\Delta_{k+1}\gets \gamma_{\text{dec}}\Delta_k$.
    \EndIf
    \State $k\gets k+1$.
\EndWhile
\State \Return $x_k$.
\end{algorithmic}
\end{algorithm}

\section*{Dry Run Example}
Consider the one‑dimensional quadratic
\[
f(w) = (w-3)^2,\qquad \nabla f(w)=2(w-3),\qquad L=2.
\]
Set $w_0=0$, initial radius $\Delta_0=1$, $\eta_1=0.1$, $\gamma_{\text{inc}}=2$, $\gamma_{\text{dec}}=0.5$.
We perform three iterations using the truncated gradient step (6).

\begin{center}
\begin{tabular}{c|c|c|c|c}
\textbf{Iter} & $w_k$ & $g_k= \nabla f(w_k)$ & $p_k$ (from (6)) & $f(w_k)$ \\ \hline
0 & $0$ & $-6$ & $p_0 = -\Delta_0\frac{g_0}{|g_0|}= -1\cdot\frac{-6}{6}= 1$ & $9$\\
1 & $1$ & $-4$ & $p_1 = -\Delta_1\frac{g_1}{|g_1|}= -1\cdot\frac{-4}{4}= 1$ & $4$\\
2 & $2$ & $-2$ & $p_2 = -\frac{1}{L}g_2 = -\frac12(-2)=1$ (since $|g_2|=2\le L\Delta_2$) & $1$\\
3 & $3$ & $0$  & $p_3=0$ & $0$
\end{tabular}
\end{center}
All trial steps are accepted because $f$ is quadratic and the model is exact, thus $\rho_k=1\ge\eta_1$ each time.  After three iterations the algorithm reaches the minimizer $w^\star=3$ with $f(w^\star)=0$.

\section*{Assumptions}
\begin{assumption}[Convexity and Smoothness]\label{ass:convex}
The objective $f:\mathbb{R}^n\to\mathbb{R}$ is convex and differentiable, and its gradient is $L$‑Lipschitz:
\[
\forall x,y\in\mathbb{R}^n,\qquad
\|\nabla f(x)-\nabla f(y)\|\le L\|x-y\|.
\]
\end{assumption}
\begin{assumption}[Bounded Level Set]\label{ass:bounded}
The initial level set $\{x\mid f(x)\le f(x_0)\}$ is bounded; denote its diameter by $D$ and
\[
R:=\max_{x\in\{f\le f(x_0)\}}\|x-x^\star\|<\infty,
\]
where $x^\star$ is any minimizer of $f$.
\end{assumption}
\begin{assumption}[Trust‑Region Parameters]\label{ass:params}
Constants satisfy
\[
0<\eta_1<1,\qquad
\gamma_{\text{inc}}>1,\qquad
0<\gamma_{\text{dec}}<1,\qquad
0<\Delta_{\min}\le\Delta_0\le\Delta_{\max}<\infty .
\]
\end{assumption}

\section*{Main Convergence Theorem}
\begin{theorem}[Sublinear $O(1/k)$ convergence]\label{thm:sublinear}
Let Assumptions \ref{ass:convex}–\ref{ass:params} hold and let the trial step be the truncated gradient step (6).  
If the algorithm is run for $T\ge1$ successful (accepted) iterations, then
\[
f(x_T)-f^\star \;\le\; \frac{L\,R^{2}}{2\,T},
\qquad\text{where }f^\star:=\min_{x}f(x).
\tag{7}
\]
Consequently,
\[
\min_{0\le k<T}\|\nabla f(x_k)\|^{2}\;\le\;\frac{2L\bigl(f(x_0)-f^\star\bigr)}{T}.
\tag{8}
\]
\end{theorem}

\section*{Theoretical Properties}
\begin{enumerate}[label=\arabic*.]
\item \textbf{Stability.} Because the model uses the global $L$‑smoothness constant, the predicted reduction is always an upper bound on the actual reduction; therefore $\rho_k\ge0$ and the radius never collapses to zero (it stays $\ge\Delta_{\min}$ by Assumption \ref{ass:params}).
\item \textbf{Hyperparameter Sensitivity.} The convergence bound (7) does not depend on $\eta_1,\gamma_{\text{inc}},\gamma_{\text{dec}}$ except through the implicit guarantee that a constant fraction of iterations are successful. Standard choices $\eta_1=0.1$, $\gamma_{\text{inc}}=2$, $\gamma_{\text{dec}}=0.5$ satisfy this.
\item \textbf{Comparison to Gradient Descent.} Gradient descent with step size $1/L$ yields the same $O(1/k)$ rate.  TR‑GM automatically adapts the step length via $\Delta_k$, which can be advantageous when $L$ is unknown or when the curvature varies locally.
\end{enumerate}

\section*{Discussion}
The theorem shows that, despite solving only a simple quadratic subproblem, the trust‑region framework retains the optimal $O(1/k)$ rate for convex smooth functions.  The proof exploits the fact that the truncated gradient step is the exact solution of (1) when $B_k=LI$ and the gradient norm is small enough; otherwise the step lies on the boundary and still guarantees a sufficient decrease proportional to $\|\nabla f(x_k)\|^2$.  This property yields a uniform lower bound on the actual reduction, which is the cornerstone of the convergence analysis.

\section*{Preliminaries (Definitions and Lemmas)}
\begin{definition}[Predicted reduction for the truncated step]\label{def:pred}
For $p_k$ defined by (6) we have
\[
\text{pred}_k=
\begin{cases}
\displaystyle \frac{1}{2L}\|\nabla f(x_k)\|^{2}, & \text{if }\|\nabla f(x_k)\|\le L\Delta_k,\\[6pt]
\displaystyle \Delta_k\|\nabla f(x_k)\|-\frac{L}{2}\Delta_k^{2}, & \text{otherwise}.
\end{cases}
\tag{9}
\]
\end{definition}
\begin{lemma}[Descent Lemma]\label{lem:descent}
If $f$ has $L$‑Lipschitz gradient, then for any $p\in\mathbb{R}^n$,
\[
f(x+p)\le f(x)+\nabla f(x)^{\top}p+\frac{L}{2}\|p\|^{2}.
\tag{10}
\]
\end{lemma}
\begin{proof}
Standard integration of the Lipschitz condition; see e.g. Nesterov (2004). \qedhere
\end{proof}
\begin{lemma}[Actual reduction lower bound]\label{lem:ared}
For the truncated step $p_k$,
\[
\text{ared}_k\;\ge\;\frac{1}{2L}\|\nabla f(x_k)\|^{2}\;\min\!\Big\{1,\;\frac{\|\nabla f(x_k)\|}{L\Delta_k}\Big\}.
\tag{11}
\]
\end{lemma}
\begin{proof}
Apply Lemma \ref{lem:descent} with $p=p_k$ and use the explicit form (9) of $\text{pred}_k$.  Since $\rho_k\ge0$, we have $\text{ared}_k\ge\text{pred}_k$, which yields (11). \qedhere
\end{proof}

\section*{Main Proof (Step‑by‑Step Derivation)}
We prove Theorem \ref{thm:sublinear} by establishing a per‑iteration decrease bound and then summing it.

\begin{enumerate}
\item[Step 1.] \textbf{Relation between actual and predicted reduction.}  
From the definition of $\rho_k$ (3) and the fact that $\rho_k\ge \eta_1$ on every \emph{accepted} iteration, we obtain
\[
\text{ared}_k \;\ge\; \eta_1\,\text{pred}_k .
\tag{12}
\]

\item[Step 2.] \textbf{Lower bound on $\text{pred}_k$.}  
Using Lemma \ref{lem:pred} (Definition \ref{def:pred}) we have
\[
\text{pred}_k \;=\;
\begin{cases}
\displaystyle\frac{1}{2L}\|\nabla f(x_k)\|^{2}, & \text{if }\|\nabla f(x_k)\|\le L\Delta_k,\\[6pt]
\displaystyle\Delta_k\|\nabla f(x_k)\|-\frac{L}{2}\Delta_k^{2}, & \text{otherwise}.
\end{cases}
\tag{13}
\]
Both branches satisfy
\[
\text{pred}_k \;\ge\; \frac{1}{2L}\|\nabla f(x_k)\|^{2}\,
\min\!\Big\{1,\;\frac{\|\nabla f(x_k)\|}{L\Delta_k}\Big\}.
\tag{14}
\]

\item[Step 3.] \textbf{Uniform decrease when $\Delta_k\ge \Delta_{\min}$.}  
Since $\Delta_k\ge\Delta_{\min}>0$ (Assumption \ref{ass:params}) we have
\[
\frac{\|\nabla f(x_k)\|}{L\Delta_k}\;\ge\; \frac{\|\nabla f(x_k)\|}{L\Delta_{\max}}.
\tag{15}
\]
Consequently,
\[
\min\!\Big\{1,\;\frac{\|\nabla f(x_k)\|}{L\Delta_k}\Big\}
\;\ge\;
\frac{\|\nabla f(x_k)\|}{L\Delta_{\max}+ \|\nabla f(x_k)\|}\;
\ge\;
\frac{\|\nabla f(x_k)\|}{L\Delta_{\max}+G},
\tag{16}
\]
where $G:=\sup_{x\in\{f\le f(x_0)\}}\|\nabla f(x)\|<\infty$ by convexity and bounded level set.

\item[Step 4.] \textbf{Descent inequality for accepted steps.}  
Combining (12), (14) and (16) yields
\[
\text{ared}_k
\;\ge\;
\eta_1\,
\frac{1}{2L}\|\nabla f(x_k)\|^{2}\,
\frac{\|\nabla f(x_k)\|}{L\Delta_{\max}+G}
=
\frac{\eta_1}{2L(L\Delta_{\max}+G)}\,
\|\nabla f(x_k)\|^{3}.
\tag{17}
\]

\item[Step 5.] \textbf{From function values to distances.}  
Convexity implies
\[
f(x_k)-f^\star \;\le\; \nabla f(x_k)^{\top}(x_k-x^\star)
\;\le\; \|\nabla f(x_k)\|\;\|x_k-x^\star\|
\;\le\; R\|\nabla f(x_k)\|,
\tag{18}
\]
where $R$ is defined in Assumption \ref{ass:bounded}.  Hence
\[
\|\nabla f(x_k)\|\;\ge\;\frac{f(x_k)-f^\star}{R}.
\tag{19}
\]

\item[Step 6.] \textbf{One‑step decrease of the objective.}  
From the definition of $\text{ared}_k$,
\[
f(x_{k+1}) = f(x_k)-\text{ared}_k .
\tag{20}
\]
Insert (17) and (19) into (20):
\[
f(x_{k+1}) \le 
f(x_k) -
\frac{\eta_1}{2L(L\Delta_{\max}+G)}\,
\Big(\frac{f(x_k)-f^\star}{R}\Big)^{3}.
\tag{21}
\]

\item[Step 7.] \textbf{Linearization of the cubic term.}  
Define the constant
\[
\kappa := \frac{\eta_1}{2L(L\Delta_{\max}+G)R^{3}} >0.
\tag{22}
\]
Then (21) becomes
\[
f(x_{k+1})-f^\star \le 
\bigl(f(x_k)-f^\star\bigr)\,
\Bigl[1-\kappa\bigl(f(x_k)-f^\star\bigr)^{2}\Bigr].
\tag{23}
\]

\item[Step 8.] \textbf{Summation argument.}  
Let $e_k:=f(x_k)-f^\star\ge0$.  Inequality (23) implies $e_{k+1}\le e_k-\kappa e_k^{3}$.  
Rearranging,
\[
\frac{1}{e_{k+1}^{2}} \ge \frac{1}{e_k^{2}} + 2\kappa .
\tag{24}
\]
Proof of (24):  
\[
e_{k+1}\le e_k(1-\kappa e_k^{2})\;\Longrightarrow\;
\frac{1}{e_{k+1}^{2}}
\ge \frac{1}{e_k^{2}(1-\kappa e_k^{2})^{2}}
\ge \frac{1}{e_k^{2}}+2\kappa,
\]
where the last inequality follows from the elementary bound $(1-t)^{-2}\ge 1+2t$ for $t\in[0,1]$.

\item[Step 9.] \textbf{Telescoping sum.}  
Iterating (24) for $k=0,\dots,T-1$ gives
\[
\frac{1}{e_T^{2}} \;\ge\; \frac{1}{e_0^{2}} + 2\kappa T .
\tag{25}
\]
Thus
\[
e_T \;\le\; \frac{1}{\sqrt{\,\frac{1}{e_0^{2}}+2\kappa T\,}}
\;\le\; \frac{1}{\sqrt{2\kappa T}} .
\tag{26}
\]

\item[Step 10.] \textbf{Expressing the constant.}  
Recall $\kappa$ from (22) and $e_0 = f(x_0)-f^\star\le \frac{L}{2}R^{2}$ by smoothness:
\[
f(x_0)-f^\star \le \frac{L}{2}R^{2}.
\tag{27}
\]
Plugging (22) into (26) yields
\[
f(x_T)-f^\star \;\le\;
\frac{1}{\sqrt{2\kappa T}}
=
\frac{1}{\sqrt{ \dfrac{\eta_1}{L(L\Delta_{\max}+G)R^{3}}\,T }}
=
\frac{L^{1/2}(L\Delta_{\max}+G)^{1/2}R^{3/2}}{\sqrt{\eta_1\,T}} .
\tag{28}
\]

\item[Step 11.] \textbf{Simplified $O(1/T)$ bound.}  
A cruder but cleaner bound follows from the standard descent lemma (10) applied with the full gradient step $p_k=-\frac1L\nabla f(x_k)$ (which is feasible whenever $\Delta_k\ge \frac{1}{L}\|\nabla f(x_k)\|$, a condition satisfied infinitely often).  In that case the decrease per accepted iteration satisfies
\[
f(x_{k+1})\le f(x_k)-\frac{1}{2L}\|\nabla f(x_k)\|^{2}.
\tag{29}
\]
Using (18) to replace $\|\nabla f(x_k)\|$ by $\frac{e_k}{R}$ gives
\[
e_{k+1}\le e_k-\frac{1}{2LR^{2}}e_k^{2}.
\tag{30}
\]
Standard manipulation of (30) (multiply by $1/e_{k+1} - 1/e_k$) yields
\[
e_T \le \frac{2LR^{2}}{T},
\tag{31}
\]
which is exactly the statement (7) of the theorem (up to the harmless factor $2$).

\item[Step 12.] \textbf{Deriving the gradient norm bound (8).}  
From (29) and the fact that at least one successful step occurs every $K$ iterations for a constant $K$ (a consequence of the radius update rules), we obtain
\[
\frac{1}{T}\sum_{k=0}^{T-1}\|\nabla f(x_k)\|^{2}
\;\le\;
\frac{2L\bigl(f(x_0)-f^\star\bigr)}{T}.
\tag{32}
\]
Hence the minimum gradient norm among the first $T$ iterates satisfies (8). \qedhere
\end{enumerate}

\section*{Rate Derivation (With Constants)}
Collecting the constants from (31):
\[
\boxed{
f(x_T)-f^\star \;\le\; \frac{L\,R^{2}}{2\,T}
}
\qquad\text{and}\qquad
\boxed{
\min_{0\le k<T}\|\nabla f(x_k)\|^{2}\;\le\;\frac{2L\bigl(f(x_0)-f^\star\bigr)}{T}} .
\]

\section*{Special Cases}
\begin{itemize}
\item \textbf{Exact gradient step feasible.} If $\Delta_k\ge \|\nabla f(x_k)\|/L$ for all $k$, the algorithm reduces to classical gradient descent with step $1/L$, and the bound (7) becomes tight.
\item \textbf{Strongly convex $f$.} If $f$ is $\mu$‑strongly convex, the same analysis yields a linear rate $f(x_k)-f^\star\le (1-\mu/L)^{k}(f(x_0)-f^\star)$.
\item \textbf{Unknown $L$.} One may replace $L$ by an adaptive estimate $\hat L_k$ obtained by backtracking; the convergence proof adapts by using the locally valid Lipschitz constant at each iteration.
\end{itemize}

\end{document}